\section{Methods}
\label{sec:methods}

\subsection{Data Source, Skip-Logic Audit, and Pre-Incision Feature Partition}
We analyzed all 368,132 isolated-CABG operations recorded in STS ACSD version~2.81 between July~1, 2014, and June~30, 2017, spanning 1,091 U.S.\ surgical centers across all nine U.S.\ Census divisions.\cite{sts2014v281,shahian2018sts,obrien2018sts2} This secondary analysis of de-identified national quality-improvement registry data was conducted under institutional Data Use Agreement governance; no patient-level identifiers or re-operation linkages were accessed, and reporting adheres strictly to TRIPOD+AI (Supplementary Table~\ref{tab:supp_tripod}) and PROBAST+AI guidelines.\cite{collins2024tripodai,moons2025probastai}

Before any model development, we performed a forensic three-part data quality and leakage audit (full ledger in Supplementary Section~S1):
\begin{enumerate}
  \item \textbf{Skip-logic outcome reconstruction:} Under STS Sequence~4930, postoperative atrial fibrillation or flutter requiring treatment (\cli{COtAFib}) is populated only when the parent complication gate \cli{Complics == 1}. Uncomplicated surgical courses (\cli{Complics == 2}; 232,783 cases, 63.31\%) leave \cli{COtAFib} structurally missing; dropping missing \cli{COtAFib} rows induces severe collider bias by conditioning on post-surgical complications.\cite{hernan2020causal} Resolving parent--child skip logic yielded \textbf{367,695 labeled procedures} (437 indeterminate records excluded) with \textbf{89,337 POAF events (24.30\%)}. External polarity checks confirmed that POAF-positive patients were 5.3~years older (68.5 vs.\ 63.2~years) and received discharge amiodarone (74.3\% vs.\ 14.7\%) and warfarin (16.5\% vs.\ 1.8\%) markedly more often.
  \item \textbf{Exhaustive 638 $\to$ 207 pre-incision variable audit:} Every column was audited individually against v2.81 data specifications rather than sliced by column index. Naive positional slicing prior to operative approach (\cli{OPApp}) inadvertently retains seven hidden post-incision or operative-duration leaks: \cli{LOS}, \cli{PLOS}, \cli{ORDuration}, \cli{SIDuration}, \cli{SurgtoMort}, \cli{SurgtoRead}, and \cli{SurgtoSternInf}. All seven fields---together with cardiopulmonary bypass and cross-clamp times, discharge therapies, STS risk-calculator outputs, auxiliary complications, and hospital/surgeon identifiers---were permanently excluded, leaving strictly \textbf{207 pre-incision predictors} (206 numeric, ordinal, or binary clinical variables plus 9-level \cli{Region}).
  \item \textbf{Sentinel decoding and preoperative AF sensitivity flag:} Non-monotonic ``Not Documented'' codes (e.g., code~5 on valve regurgitation, codes~5--6 on chronic lung disease) were decoded prior to clinical scoring. Under v2.81 coding, 1,900 cases (0.52\%; 138 in external validation) carried recorded paroxysmal or persistent preoperative AF (\cli{ArrhythAFib} codes~2--3); we retained an explicit sensitivity flag and report primary performance on both the full cohort and the \emph{no recorded preoperative AF} subgroup ($n=21{,}103$).
\end{enumerate}

\subsection{Forensic Critique of Naive 2017 Splits and Dual Evaluation Regimes}
We audited five partition geometries across all 367,695 labeled cases (Figure~\ref{fig:study_design}B--C; Supplementary Section~\ref{supp:schema}). The conventional temporal split---training on 2014--2016 ($N=306{,}143$) and testing on calendar year 2017 ($N=61{,}552$)---fails both as an external-hospital test and as a stationary statistical benchmark for five reasons: (1)~99.5\% hospital overlap (998 of 1,003 hospitals operating in 2017 appear in training; only 5 centers are unseen); (2)~threefold inter-hospital POAF dispersion (12.2\%--34.2\% across 928 centers with $\ge 100$ cases, 5th--95th percentile; overall 11.7\%--34.8\%, SD~6.68--6.97~pp); (3)~half-year boundary asymmetry (v2.81 spans July~2014--June~2017, making 2014 H2-only at 23.69\% POAF vs.\ 2017 H1-only at 24.83\%); (4)~non-stationary missingness drift ($+6.05$~pp on \cli{LVSD}, $+6.04$~pp on \cli{LVEDD}, $+5.60$~pp on \cli{PctStenAM}; $\mathrm{Max}\,|\mathrm{SMD}| = 0.2793$); and (5)~substantial multivariate covariate shift (adversarial classifier AUC 0.6208 overall; 0.5908 hospital-grouped; 0.7661 vs.\ 2014 alone).

To cleanly separate external transportability from in-network institutional deployment, we constructed two rigorously matched primary evaluation regimes using identical sample sizes ($N=21{,}241$ natural-prior holdout; nested $N=50{,}000$ class-balanced development benchmark for frontier tabular foundation models\cite{jaeger2026tabpfn35,qu2026kumotabular,kong2026tabfm}):

\runin{Regime 1: Cold-Start Unseen-Hospital External Validation (0.0\% Hospital Overlap)}
Random hospital-grouped splitting removes site overlap but produces severe finite-cluster imbalance ($\mathrm{Max}\,|\mathrm{SMD}| = 0.3572$, Region TV $= 0.1645$, POAF gap 1.13~pp). We therefore applied cluster-level multivariate anticlustering\cite{papenberg2021anticlustering} minimizing covariate proxy loss plus $50\times$ 32-dimensional whitened SPlit Energy Distance\cite{joseph2022split,szekely2013energy} (monitoring multi-kernel $\mathrm{MMD}^2$\cite{gretton2012kernel}) subject to a $\ge 5{,}000$ holdout event floor. This partitioned the 1,091 centers into \textbf{1,017 development hospitals} ($N=346{,}454$; 94.2\%) and \textbf{74 strictly unseen validation hospitals} ($N=21{,}241$; 5.8\%, 5,160 events, 24.29\%, 0.0\% hospital overlap), cutting Energy Distance 4.1$\times$ ($0.01732 \to 0.00427$) and $\mathrm{Max}\,|\mathrm{SMD}|$ 3.5$\times$ ($0.3572 \to 0.1028$). Inside the exact same 1,017 development hospitals, class-conditional support-point subsampling\cite{joseph2022split} created a balanced 50,000-patient benchmark (25,000 POAF / 25,000 controls across five exact 10,000-row hospital-disjoint CV folds of 203--204 centers each).

\runin{Regime 2: In-Network Seen-Hospital Deployment (100.0\% All-Hospitals-In-Train)}
To directly quantify the effect of hospital familiarity free of calendar confounding and mirror deployment within an established national registry or health system, we constructed a matched \textbf{100.0\% Seen-Hospital} partition across all 367,695 cases in which \emph{every single one of the 1,091 U.S.\ surgical centers} appears in both the full development set ($N=346{,}454$; 1,091/1,091 hospitals) and the nested 50,000-patient balanced training set (1,091/1,091 hospitals; 5 stratified 10,000-row folds), with zero patient overlap and an exact matched natural-prior holdout of $N=21{,}241$ patients across 1,086 seen hospitals (0 unseen hospitals; 24.30\% POAF). On this seen-hospital cohort, we evaluated every algorithm under both (a)~\textbf{patient-only 207 raw predictors} (isolating implicit hospital familiarity) and (b)~\textbf{in-network site-calibrated deployment}, incorporating each hospital's strictly train-only Empirical-Bayes (EB) shrunk historical log-odds deviation $\hat{\delta}_h = \operatorname{logit}\!\big((e_h + m_0\pi)/(n_h + m_0)\big) - \operatorname{logit}(\pi)$ (computed strictly out-of-fold inside training folds; zero test-label contact), both via parameter-free closed-form hierarchical Bayes site-intercept updating ($\operatorname{Platt}_{\mathrm{OOF}} + \text{dev offset} + \hat{\delta}_h$) and native 208-feature retraining.

\subsection{Comparators, Evolutionary Continuous-Target Learning, and Parallel Genetic Programming}
Every 50,000-patient benchmark model consumed strictly the same 207 raw pre-incision variables under a process-wide sealed-label audit hook blocking any test-label access during training or inference.

\runin{Clinical Scores and Six Binary Tabular ELO Baselines}
We evaluated six established clinical comparators calibrated strictly on the 50k development cohort (Mariscalco POAF,\cite{mariscalco2014poaf} Burgos COM-AF,\cite{burgos2021comaf} $\mathrm{CHA_2DS_2\text{-}VASc}$,\cite{lip2010cha2ds2vasc} HATCH,\cite{devos2010hatch} preoperative Mathew JAMA proxy,\cite{mathew2004multicenter} and continuous age alone) alongside six world-leading binary tabular model families:\cite{grinsztajn2022why,erickson2025tabarena} tuned 5-fold \textbf{CatBoost},\cite{prokhorenkova2018catboost} \textbf{LightGBM},\cite{ke2017lightgbm} and \textbf{XGBoost},\cite{chen2016xgboost} plus three frontier zero-shot in-context tabular foundation models---\textbf{PriorLabs TabPFN-3.5} (native 50k context and $5\times 10\text{k}$ bagged),\cite{hollmann2025tabpfn,jaeger2026tabpfn35} \textbf{NVIDIA KumoTabular-Large v2} ($5\times 10\text{k}$ with mandatory 256-row train-only class-orientation guard),\cite{qu2026kumotabular} and \textbf{Google TabFM v1.0.0} ($5\times 10\text{k}$, 80 NNLS estimators).\cite{kong2026tabfm} Full-registry 346,454-case CatBoost served as the 7$\times$ data-scaling reference, and 92 agent-engineered clinical composites were audited independently across both cohort scales (Supplementary Section~S2).

\runin{EvoTarget-FT: 1,000-Generation Target Evolution and Strictly Nested 2$\times$2 Causal Audit}
In Stage~1, an evolutionary particle-swarm solver distills\cite{hinton2015distilling} continuous training targets $z^{(g)} \in [0,1]^{50{,}000}$ strictly inside the development cohort across whole-hospital shuffled multi-family teachers (CatBoost, LightGBM, XGBoost, and L2 Ridge; 256 fits per generation) using exponentially-weighted moving-average (EWMA) step hardening $z^{(g+1)}=(1-\beta_g)z^{(g)}+\beta_g\tilde{z}^{(g)}$ while gene-locking the 20\% most confident consensus cases (full equations in Supplementary Section~S3). We executed 1,000 full generations (256,000 fits) to audit late-stage stability and selected Generation~20 ($z^*$) for Stage~2 training of \textbf{EvoTarget-FT}, a Dual-View Periodic Linear-ReLU (PLR) clinical self-attention transformer\cite{gorishniy2021revisiting,gorishniy2022embeddings} (4 seeds $\times$ 5 folds, top-3 validation SWA,\cite{izmailov2018swa} composite cross-entropy plus pairwise ranking loss) consuming strictly 207 raw predictors. To resolve any concern regarding global outer-fold target search or 289$\to$207 teacher-feature asymmetry, we also executed a \textbf{strictly nested 5-fold outer-loop} Gen-20 evolution using strictly 207 raw predictors inside each fold's $\sim$814 training hospitals alongside a global 207-raw control, forming an exact 2$\times$2 causal decomposition.

\runin{JAIN Ultra GP v2 and Dual-Paradigm Fusion}
Genetic programming synthesizes explicit symbolic feature pipelines paired with linear, spline, kernel, and tree heads. Profiling default legacy JAIN Ultra uncovered six engineering bottlenecks (single-worker search, 1~GiB worker memory cap, 20,000-row fit cap, random IID internal splits leaking hospital incidence, RMSE fitness on smooth targets, and single-model refits). Our re-engineered build (\textbf{JAIN Ultra GP v2}) resolves all six via 24 parallel GP workers, 8~GiB memory, full 50k training, 64-bin Newton soft-BCE boosting (\cli{histgbm2}) with optional binary-$y$ early stopping (\cli{with\_stop\_target}), hospital-grouped CV optimizing binary ROC-AUC, CUDA learned-ranker guidance, and 64-thread fold-bagging across six evolved operator families (\cli{ensemble}, \cli{winsor\_clip} $\to$ \cli{pair\_ratio} $\to$ \cli{sigmoid} $\to$ \cli{ridge\_gcv}, \cli{pair\_ratio} $\to$ \cli{elasticnet}, \cli{histgbm2}, \cli{oblique\_hinge} $+$ \cli{cheb} $+$ \cli{pair\_prod} $\to$ \cli{ridge}, and \cli{histgbm}). Finally, \textbf{Dual-Paradigm Fusion} blends OOF-weighted log-odds from JAIN Ultra GP v2 and EvoTarget-FT strictly via development OOF optimization.

\subsection{Calibration, Statistical Inference, Clinical Utility, and Equity Analysis}
All 50,000-patient models trained on balanced 50/50 data were mapped back to natural clinical prevalence using uniform fold-OOF Platt scaling\cite{platt1999probabilistic} plus the closed-form strictly development-only Bayes prior-shift offset:\cite{king2001logistic,saerens2002adjusting}
\begin{equation}
  \operatorname{logit}\!\bigl(\hat{p}_{\mathrm{nat}}(x)\bigr) = \operatorname{logit}\!\bigl(\hat{p}_{\mathrm{bal}}(x)\bigr) - 1.136479,
  \label{eq:bayes_shift}
\end{equation}
where $1.136479 = \operatorname{logit}(0.50) - \operatorname{logit}(0.242967)$ uses strictly the development full-training prevalence ($-1.136495$ in the seen-hospital split). Calibration was assessed via 15-bin Expected Calibration Error (ECE),\cite{guo2017calibration,vancalster2019calibration} Integrated Calibration Index (ICI),\cite{austin2019ici} logistic calibration slope, calibration-in-the-large (O:E ratio), Brier score, and decile risk plots.

Because patients cluster within hospitals, all 95\% confidence intervals and paired model comparisons used \textbf{5,000-draw whole-hospital cluster bootstrapping}\cite{davison1997bootstrap,obuchowski1997clustered} (seed~42) resampling all 74 unseen hospitals (Regime~1) or all 1,086 seen hospitals (Regime~2) as intact clusters. Clinical utility was quantified via Decision Curve Analysis (Net Benefit and net unnecessary prophylactic interventions avoided per 100 patients across threshold probabilities 10\%--45\%).\cite{vickers2006decision} Equity and transportability were audited across 25 prespecified demographic, clinical, acuity, and hospital-volume strata, 74 individual hospital funnels, and chronological prospective site-onboarding recalibration (2014--2016 $\to$ 2017). Clinical face validity was evaluated via standardized multivariable surrogate regression ($R^2$ on champion log-odds).
